% GENERATED FILE. Edit canonical lessons, then run npm run generate:books.
% canonical-source-hash: fnv1a64:0a6c8556ecd6a024
% canonical-lessons: ZH-C391-richeng, ZH-C391-linshi, ZH-C391-laideji, ZH-C391-ganshang, ZH-C391-daoshihou, ZH-R391-review-386-391

\chapter{When Plans Change}
\label{ch:zh-when-plans-change}
\hlchaptermodality{391}

\begin{chapteropening}
\textbf{By the end of this chapter:} \emph{I can talk about a day's schedule, things at the last minute, still having time, catching a train and what we will do when the time comes.}

Complete the last lesson of chapter 391: I can talk about a day's schedule, things at the last minute, still having time, catching a train and what we will do when the time comes.
\end{chapteropening}

\section[rìchéng]{rìchéng — a day's schedule}

\label{lesson:ZH-C391-richeng}



\begin{quote}
Before the new one: say the Mandarin for on the agreed date, then the Mandarin for approximately.
\end{quote}

\subsection*{You'll want to know: rìchéng}
\emph{rìchéng} — \textquotedblleft{}a day's schedule\textquotedblright{}.

\begin{quote}
\emph{míngtiān de rìchéng hěn mǎn} — Tomorrow's schedule is full.
\end{quote}

Its characters come later; for now it is a word for the ear.

\begin{cousinweb}
\emph{rìchéng} is built from \zh{日} 'a day' and \emph{chéng} 'a stage of a journey', both inherited from Middle Chinese, \emph{nyit} and \emph{drjeng} (Guangyun, in Baxter's transcription), which glosses the second 'a set term' and 'a limit'.
\end{cousinweb}

\begin{grammarlens}[title={What you've built}]
One more word for saying what would happen.
\end{grammarlens}

\subsection*{Your turn}
\begin{itemize}
\raggedright
  \item \emph{Say it:} \emph{rìchéng}
  \item \emph{Say it:} \emph{rìchéng}, once more
  \item \emph{Say it:} \emph{míngtiān de rìchéng hěn mǎn} — Tomorrow's schedule is full.
  \item \emph{Recall:} say the Mandarin for on the agreed date, then the Mandarin for approximately, then say \emph{rìchéng} again
\end{itemize}

\begin{tcolorbox}[breakable,colback=teal!4,colframe=teal!35!black,title={Before you move on}]
What does rìchéng mean? (\textquotedblleft{}A day's schedule\textquotedblright{}.) Say it once more.
\end{tcolorbox}
\section[línshí]{línshí — at the last minute; for the time being}

\label{lesson:ZH-C391-linshi}



\begin{quote}
Before the new one: say the Mandarin for approximately, then the Mandarin for a day's schedule.
\end{quote}

\subsection*{You'll want to know: línshí}
\emph{línshí} — \textquotedblleft{}at the last minute; for the time being\textquotedblright{}.

\begin{quote}
\emph{wǒ línshí yǒushì bù néng qù le} — Something's come up at the last minute and I can't go.
\end{quote}

Its characters come later; for now it is a word for the ear.

\begin{cousinweb}
\emph{línshí} is built from \emph{lín} 'to come close' and \zh{时} 'the time', both inherited from Middle Chinese, \emph{lim} and \emph{dzyi} (Guangyun, in Baxter's transcription), which glosses the first 'to be present at': when the moment is upon you.
\end{cousinweb}

\begin{grammarlens}[title={What you've built}]
One more word for saying what would happen.
\end{grammarlens}

\subsection*{Your turn}
\begin{itemize}
\raggedright
  \item \emph{Say it:} \emph{línshí}
  \item \emph{Say it:} \emph{línshí}, once more
  \item \emph{Say it:} \emph{wǒ línshí yǒushì bù néng qù le} — Something's come up at the last minute and I can't go.
  \item \emph{Recall:} say the Mandarin for approximately, then the Mandarin for a day's schedule, then say \emph{línshí} again
\end{itemize}

\begin{tcolorbox}[breakable,colback=teal!4,colframe=teal!35!black,title={Before you move on}]
What does línshí mean? (\textquotedblleft{}At the last minute; for the time being\textquotedblright{}.) Say it once more.
\end{tcolorbox}
\section[láidejí]{láidejí — to still have time (to do something)}

\label{lesson:ZH-C391-laideji}



\begin{quote}
Before the new one: say the Mandarin for a day's schedule, then the Mandarin for at the last minute.
\end{quote}

\subsection*{You'll want to know: láidejí}
\emph{láidejí} — \textquotedblleft{}to still have time (to do something)\textquotedblright{}.

\begin{quote}
\emph{rúguǒ xiànzài zǒu hái láidejí} — If we leave now, there's still time.
\end{quote}

Its characters come later; for now it is a word for the ear.

\begin{cousinweb}
\emph{láidejí} is built from \zh{来} 'to come', \zh{得} 'can' and \emph{jí} 'to reach', all inherited from Middle Chinese, \emph{loj}, \emph{tok} and \emph{gip} (Guangyun, in Baxter's transcription), which glosses the last 'to get to': you can still get there. It is the other side of láibují.
\end{cousinweb}

\begin{grammarlens}[title={What you've built}]
One more word for saying what would happen.
\end{grammarlens}

\subsection*{Your turn}
\begin{itemize}
\raggedright
  \item \emph{Say it:} \emph{láidejí}
  \item \emph{Say it:} \emph{láidejí}, once more
  \item \emph{Say it:} \emph{rúguǒ xiànzài zǒu hái láidejí} — If we leave now, there's still time.
  \item \emph{Recall:} say the Mandarin for a day's schedule, then the Mandarin for at the last minute, then say \emph{láidejí} again
\end{itemize}

\begin{tcolorbox}[breakable,colback=teal!4,colframe=teal!35!black,title={Before you move on}]
What does láidejí mean? (\textquotedblleft{}To still have time (to do something)\textquotedblright{}.) Say it once more.
\end{tcolorbox}
\section[gǎnshàng]{gǎnshàng — to catch in time; to make it onto a train}

\label{lesson:ZH-C391-ganshang}



\begin{quote}
Before the new one: say the Mandarin for at the last minute, then the Mandarin for to still have time.
\end{quote}

\subsection*{You'll want to know: gǎnshàng}
\emph{gǎnshàng} — \textquotedblleft{}to catch in time; to make it onto a train\textquotedblright{}.

\begin{quote}
\emph{yàoshi wǒmen kuài yìdiǎnr jiù néng gǎnshàng huǒchē} — If we're a bit quicker, we can catch the train.
\end{quote}

Its characters come later; for now it is a word for the ear.

\begin{cousinweb}
\emph{gǎnshàng} is built from \emph{gǎn} 'to chase' and \zh{上} 'to go up'. The second is inherited from Middle Chinese \emph{dzyangX} (Guangyun, in Baxter's transcription), glossed 'to climb'; the dictionary has no entry for the first, so its history is not settled.
\end{cousinweb}

\begin{grammarlens}[title={What you've built}]
One more word for saying what would happen.
\end{grammarlens}

\subsection*{Your turn}
\begin{itemize}
\raggedright
  \item \emph{Say it:} \emph{gǎnshàng}
  \item \emph{Say it:} \emph{gǎnshàng}, once more
  \item \emph{Say it:} \emph{yàoshi wǒmen kuài yìdiǎnr jiù néng gǎnshàng huǒchē} — If we're a bit quicker, we can catch the train.
  \item \emph{Recall:} say the Mandarin for at the last minute, then the Mandarin for to still have time, then say \emph{gǎnshàng} again
\end{itemize}

\begin{tcolorbox}[breakable,colback=teal!4,colframe=teal!35!black,title={Before you move on}]
What does gǎnshàng mean? (\textquotedblleft{}To catch in time; to make it onto a train\textquotedblright{}.) Say it once more.
\end{tcolorbox}
\section[\texorpdfstring{\zh{到时候} (dào shíhou)}{到时候 (dào shíhou)}]{\zh{到时候} (dào shíhou) — when the time comes}

\label{lesson:ZH-C391-daoshihou}



\begin{quote}
Before the new one: say the Mandarin for to still have time, then the Mandarin for to make it onto a train.
\end{quote}

\subsection*{You'll want to know: \zh{到时候}}
\textbf{\zh{到时候}} — \emph{dào shíhou} — \textquotedblleft{}when the time comes\textquotedblright{}.

\begin{quote}
\textbf{\zh{到时候我一定来}} \emph{dào shíhou wǒ yídìng lái} — When the time comes, I'll be there for sure.
\end{quote}

\begin{cousinweb}
\textbf{\zh{到时候}} (\emph{dào shíhou}) is built from \zh{到} 'to arrive' and \zh{时候} 'the time', all inherited from Middle Chinese, \emph{tawH}, \emph{dzyi} and \emph{huwH} (Guangyun, in Baxter's transcription), which glosses the last 'to wait for'.
\end{cousinweb}

\begin{grammarlens}[title={What you've built}]
One more word for saying what would happen.
\end{grammarlens}

\subsection*{Your turn}
\begin{itemize}
\raggedright
  \item \emph{Say it:} \emph{dào shíhou}
  \item \emph{Say it:} \emph{dào shíhou}, once more
  \item \emph{Say it:} \emph{dào shíhou wǒ yídìng lái} — When the time comes, I'll be there for sure.
  \item \emph{Recall:} say the Mandarin for to still have time, then the Mandarin for to make it onto a train, then say \emph{dào shíhou} again
\end{itemize}

\begin{tcolorbox}[breakable,colback=teal!4,colframe=teal!35!black,title={Before you move on}]
What does \zh{到时候} mean? (\textquotedblleft{}When the time comes\textquotedblright{}.) Say it once more.
\end{tcolorbox}
\section[(dialogue)]{Looking back: from Second-Hand to When Plans Change}

\label{lesson:ZH-R391-review-386-391}



\begin{quote}
Say the Mandarin for second-hand and to turn into.
\end{quote}

\subsection*{Your turn}
\begin{itemize}
\raggedright
  \item \emph{Say it:} second-hand, to turn into, sitting unused at home, good value for money, to turn bad (of the weather)
  \item \emph{Say it:} income, savings (in the bank), a bill (to pay), to shop online, a budget
  \item \emph{Say it:} to find time (for something), a romantic date, to miss an appointment, to do ahead of time, a time zone
  \item \emph{Say it:} spare (kept for when needed), to economise, a risk, as a precaution, to fall due
  \item \emph{Say it:} to be certain, to agree on a time, at regular intervals, on the agreed date, approximately
  \item \emph{Say it:} a day's schedule, at the last minute, to still have time, to make it onto a train, when the time comes
  \item \emph{Read it:} \zh{变} \zh{约} \zh{定} — each one aloud, then a word that holds it
\end{itemize}

\begin{tcolorbox}[breakable,colback=teal!4,colframe=teal!35!black,title={Before you move on}]
Second-hand? (\textbf{\zh{二手}}, \emph{èrshǒu}.) To turn into? (\textbf{\zh{变成}}, \emph{biànchéng}.) Sitting unused at home? (\emph{xiánzhì}.) Good value for money? (\emph{huásuàn}.) To turn bad (of the weather)? (\textbf{\zh{变天}}, \emph{biàntiān}.) Income? (\emph{shōurù}.) Savings (in the bank)? (\emph{cúnkuǎn}.) A bill (to pay)? (\emph{zhàngdān}.) To shop online? (\emph{wǎnggòu}.) A budget? (\emph{yùsuàn}.) To find time (for something)? (\emph{chōukòng}.) A romantic date? (\textbf{\zh{约会}}, \emph{yuēhuì}.) To miss an appointment? (\emph{shīyuē}.) To do ahead of time? (\emph{tíqián}.) A time zone? (\textbf{\zh{时区}}, \emph{shíqū}.) Spare (kept for when needed)? (\emph{bèiyòng}.) To economise? (\textbf{\zh{节约}}, \emph{jiéyuē}.) A risk? (\emph{fēngxiǎn}.) As a precaution? (\emph{yǐfáng}.) To fall due? (\textbf{\zh{到期}}, \emph{dàoqī}.) To be certain? (\emph{quèdìng}.) To agree on a time? (\textbf{\zh{约定}}, \emph{yuēdìng}.) At regular intervals? (\textbf{\zh{定期}}, \emph{dìngqī}.) On the agreed date? (\textbf{\zh{如期}}, \emph{rúqī}.) Approximately? (\textbf{\zh{大约}}, \emph{dàyuē}.) A day's schedule? (\emph{rìchéng}.) At the last minute? (\emph{línshí}.) To still have time? (\emph{láidejí}.) To make it onto a train? (\emph{gǎnshàng}.) When the time comes? (\textbf{\zh{到时候}}, \emph{dào shíhou}.)
\end{tcolorbox}
